% applications/examples.tex

\chapter{Application Examples}

\label{chap:applications-examples}

This chapter illustrates representative application patterns that use
established AIM concepts.

% ============================================================
\section{Example Pattern A: Transition Comparison}
% ============================================================

Two transition laws may satisfy identical boundary conditions while
producing different internal curvature evolution and therefore different
runtime-derived dynamic quantities.

A canonical indicator remains
\[
a_{\mathrm{lat}}(s)=v^2\kappa(s),
\]
showing that operational consequence depends on both geometry and speed
context.

% ============================================================
\section{Example Pattern B: Cant-Coupled Interpretation}
% ============================================================

Application interpretation uses coupled curvature/cant quantities and
runtime operators to evaluate admissibility and comfort consequences.
Derived quantities remain downstream runtime results.

% ============================================================
\section{Example Pattern C: Realized-State Comparison}
% ============================================================

Target and realized states are compared in realization-aware workflows.
The example focus is engineering consequence (maintenance or operational
adjustment), not redefinition of upstream semantics.

% ============================================================
\section{Connection to AIM}
% ============================================================

\begin{summary}
The examples demonstrate how canonical AIM concepts are used for
comparison, evaluation, and decision support in engineering contexts.

They intentionally illuminate established architecture rather than
introducing new conceptual layers.
\end{summary}
